\NumberedItem{3}
\HistoricTitle{\textsc{Gauss} an \textsc{Schumacher}.}
\HistoricCitation{[Briefwechsel zwischen \textsc{Gauss} und \textsc{Schumacher} VI, Altona 1865, S. 53.]}
\indent Ich kann nicht unterlassen, eine Unrichtigkeit, die sich in den Schluss meines letzten Briefes eingeschlichen hat, sogleich zu berichtigen. Ich schrieb diesen Schluss bloss nach dem, was ich von \textsc{Jacobis} Aufsatz (seit März 1849,
\clearpage\SourceLeaf{439}\ArithmeticHead{432}{L}{ANALYSIS. BRIEFWECHSEL.}
\noindent nach damals nur flüchtiger Ansicht) im Gedächtniss hatte oder zu haben glaubte, indem ich denselben gar nicht selbst wieder nachsah.

\indent Ich glaubte nemlich in \textsc{Jacobis} Aufsatz stehe, dass \textsc{Carlini} die Convergenz für die Mittelpunktsgleichung richtig, aber für den Radius Vector falsch angegeben habe; auch die Formel selbst hatte ich nur ihrer Form nach im Gedächtniss und meinte, dass sie mit meiner vor 40 oder mehreren Jahren [gefundenen] übereinstimmend gewesen sei. Das Wahre ist, dass meine Convergenzformel[\textsuperscript{*})] mit der von \textsc{Jacobi}[\textsuperscript{**})] übereinstimmt, nemlich wenn \(\varepsilon\) die Excentricität, \(e\) die Basis der hyp[erbolischen] Log[arithmen] bedeutet, so convergiren die Coefficienten jener Reihe langsamer als jede fallende geometrische Progression, deren Exponent kleiner ist als
\[
\frac{\varepsilon e^{\surd(1-\varepsilon\varepsilon)}}{1+\surd(1-\varepsilon\varepsilon)}
\]
\noindent oder als
\[
\operatorname{tg}\tfrac12\varphi\,{.}\,e^{\cos\varphi},\quad (\text{wenn }\varepsilon=\sin\varphi)
\]
\noindent aber etwas schneller als die geometr[ische] Progr[ession], deren Exponent dieser Grösse gleich ist.

\indent Geirrt habe ich also mich gestern ohne Zweifel\textsuperscript{***}), indem ich sagte, dass ich 1817 \textsc{Carlinis} Abhandl[ung] nur bis dahin angesehen habe, wo die Übereinstimmung in der Hauptsache (s. oben) hervorgetreten sei. Ohne Zweifel habe ich sie damals \textls[120]{gar nicht} näher angesehen, weil ich keine Lust hatte, eine 48 Seiten lange Abhandlung durchzulesen, die durch höchst verwickelte Rechnung eine Aufgabe auflösen sollte, die ich selbst (so weit es nöthig) schon lange vorher auf einer halben Octavseite aufgelöset hatte.

\indent Ich glaube mich übrigens bestimmt zu erinnern, dass ich damals, als ich meine Auflösung gefunden hatte (ich meine, in einem der ersten Jahre dieses Jahrhunderts[\textsuperscript{\ensuremath{\dagger}})], ich sogar die Richtigkeit jener Convergenzbestimmung durch
\par\medskip\noindent\rule{.24\linewidth}{.4pt}\par\smallskip
{\small
\noindent\hspace*{\parindent}[\textsuperscript{*}) Siehe oben Abschnitt [VII.], art. [2.], Gleichung [28], S. 423.]
\par\noindent\hspace*{\parindent}[\textsuperscript{**}) Siehe \textsc{Jacobi}, a. a. O., \textsc{Jacobis} Werke VII, S. 178 und 188, vergl. ferner die bald nach Abfassung dieses Briefes erschienene Bearbeitung der \textsc{Carlini}schen Abhandlung von \textsc{Jacobi}: \textit{Untersuchungen über die Convergenz der Reihe, durch welche das Keplersche Problem gelöst wird}, Astronomische Nachrichten \textbf{30}, Nr. 709—712, 1850, S. 197, \textsc{Jacobis} Werke VII, S. 189, insbesondere S. 237.]
\par}
\noindent\hspace*{\parindent}\textsuperscript{***}) Ohne Zweifel sage ich, weil ich in \textls[120]{diesem Augenblick} \textsc{Carlinis} Abhandlung nicht selbst nachgesehen habe, sondern bloss flüchtig den \textsc{Jacobi}schen Bericht.
\par{\small\noindent\hspace*{\parindent}[\textsuperscript{\ensuremath{\dagger}}) Vergl. die Angabe in dem folgenden Briefe vom 5. Februar 1850.]\par}
\clearpage\SourceLeaf{440}\ArithmeticHead{433}{R}{GAUSS AN SCHUMACHER.}
\noindent einen Fall in Concreto bei einer grossen Excentricität constatirt habe, worüber sich vielleicht noch ein Papier wird auffinden lassen.

\indent So viel heute in Eile, weil es mir unangenehm war, Ihnen gestern etwas unrichtiges aus dem Gedächtniss geschrieben zu haben.
\begin{flushright}Der Ihrige\hspace*{.15\linewidth}\end{flushright}
\noindent\hspace*{\parindent}Göttingen den 6. December 1849.\hfill C. F. \textsc{Gauss}.
\NumberedItem{4}
\HistoricTitle{\textsc{Gauss} an \textsc{Schumacher}.\quad Göttingen, 5. Februar 1850.}
\HistoricCitation{[Briefwechsel zwischen \textsc{Gauss} und \textsc{Schumacher} VI, Altona 1865, S. 58.]}
\indent Von meiner Methode, den Grad der Convergenz der nach den Cosinus und Sinus der Vielfachen eines Winkels fortschreitenden, eine beliebige periodische Function ausdrückenden Reihe zu bestimmen, habe ich noch eine numerische Rechnung, welche sich auf ein Beispiel der Mittelpunktsgleichung bezieht, aufgefunden, welches Blatt wohl \(50\pm\) Jahre alt sein mag[\textsuperscript{*})]. Die Methode leistet aber viel mehr, als bloss einen genäherten Ausdruck für ein sehr weit vom Anfange entferntes Glied zu finden; sie ist auch geeignet, alle Glieder bis zum Anfang selbst hin, numerisch zu berechnen, und zwar mit aller zu wünschenden Schärfe. In dem Maasse ist jenes Beispiel damals nicht durchgeführt, was jetzt zu ergänzen mir die Zeit fehlt. Bei weitem mehr Zeit wird aber erfordert werden, um die ganze Theorie in einer mir selbst genügenden Gestalt\textsuperscript{**}) auszuführen. Ich bin nicht abgeneigt, eine mir zu Theil werdende Musse dazu zu verwenden, möglicherweise wird aber die Arbeit dann einen grössern Umfang erhalten, als sich für Aufnahme in die A[stronomischen] N[achrichten] eignet. \(\,{.}\quad{.}\quad{.}\quad{.}\quad{.}\quad{.}\quad{.}\quad{.}\quad{.}\quad{.}\quad{.}\quad{.}\quad{.}\quad{.}\quad{.}\quad{.}\)
\par\medskip\noindent\rule{.24\linewidth}{.4pt}\par\smallskip
{\small\noindent\hspace*{\parindent}[\textsuperscript{*}) Siehe den Abschnitt VII, S. 420, insbesondere die Zahlentafel des art. 3., S. 424.]\par}
\noindent\hspace*{\parindent}\textsuperscript{**}) Sie sind ganz im Irrthum, wenn Sie glauben, dass ich darunter nur die letzte Politur in Beziehung auf Sprache und Eleganz der Darstellung verstehe. Diese kosten vergleichungsweise nur unbedeutenden Zeitaufwand; was ich meine, ist die \textls[120]{innere} Vollkommenheit. In manchen meiner Arbeiten sind solche Incidenzpunkte, die mich jahrelanges Nachdenken gekostet haben, und deren in kleinem Raum concentrirter Darstellung nachher niemand die Schwierigkeit anmerkt, die erst überwunden werden musste.
\par\smallskip\noindent{\small X1.}\hfill{\small 55}
